% Source: physical PDF pp. 171--178 (printed pp. 148--155), from the Section
% 27.6 heading through the final paragraph immediately before Section 27.7.
% Inventory: Eqs. (27.6.1)--(27.6.21), no unnumbered display groups, linked
% citations [5], [6], [7], and [3], no footnotes, Figure 27.1, one centered
% three-asterisk divider, and no tables.
% Source notation preserved: Eq. (27.6.16) uses X^* and Y^*, although
% Eqs. (27.6.5) and (27.6.13) use daggers for the corresponding adjoint
% superfields. Corrected the extra factor 1/2 in Eq. (27.6.4).
% Uncertainties: none.

\subsection{Perturbative Non-Renormalization Theorems}

From the beginning, several of the ultraviolet divergences in ordinary
renormalizable quantum field theories were found to be absent in the
supersymmetric versions of these theories. With the development in 1975 of
supergraph techniques, in which all particles in each supermultiplet are
considered together, it became possible to show that some radiative
corrections are not only finite, but are absent altogether in perturbation
theory~\hyperref[ch27-ref-5]{[5]}. Supergraphs will be described in detail in
Chapter 30, but as it happens they are not needed to prove the most important
non-renormalization theorems. This section will give a version of a method
developed by Seiberg~\hyperref[ch27-ref-6]{[6]} in 1993, which showed how the
non-renormalization theorems may be easily obtained from simple
considerations of symmetry and analyticity.

Consider a general renormalizable supersymmetric gauge theory with a number
of left-chiral superfields \(\Phi_n\) and/or gauge superfields \(V_A\). As
mentioned in Section 27.3, if we remove the factor \(g\) in the \(t_A\) and
\(C_{ABC}\) and include it instead in the gauge superfields, then the
Lagrangian density will be of the form
\begin{equation}
  \mathcal{L}
  =
  \left[\Phi^\dagger e^{-V}\Phi\right]_D
  +2\operatorname{Re}\left[f(\Phi)\right]_{\mathcal{F}}
  +\frac{1}{2g^2}\operatorname{Re}
  \left[
    \sum_{A\alpha\beta}
    \epsilon_{\alpha\beta}W_{A\alpha L}W_{A\beta L}
  \right]_{\mathcal{F}}\,,
  \tag{27.6.1}\label{eq:27.6.1}
\end{equation}
where the superpotential \(f(\Phi)\) is a gauge-invariant cubic polynomial in
the left-chiral superfields. (We are ignoring a possible \(\theta\)-term,
which has no effect in perturbation theory.)

Suppose we impose an ultraviolet cut-off \(\lambda\) on the momenta
circulating in loop graphs. As discussed in Section 12.4, we can find a
\emph{local} `Wilsonian' effective Lagrangian density \(\mathcal{L}_\lambda\)
that, with this cut-off, gives precisely the same results as the original
Lagrangian density for \(S\)-matrix elements of processes at momenta below
\(\lambda\). The effective Lagrangian density has masses and coupling
parameters that now depend on \(\lambda\), and usually there will be an
infinite number of coupling terms in the effective Lagrangian density, all
possible terms allowed by the symmetries of the theory. But things are much
simpler in supersymmetric theories. The non-renormalization theorems tell us
that, as long as the cut-off preserves supersymmetry and gauge invariance, to
all orders in perturbation theory the effective Lagrangian will have the
structure
\begin{equation}
  \begin{aligned}
  \mathcal{L}_\lambda
  ={}&
  \left[
    \mathcal{A}_\lambda
    \bigl(\Phi,\Phi^\dagger,V,\mathcal{D}\cdots\bigr)
  \right]_D
  +2\operatorname{Re}\left[f(\Phi)\right]_{\mathcal{F}}
  \\
  &+\frac{1}{2g_\lambda^2}\operatorname{Re}
  \left[
    \sum_{A\alpha\beta}
    \epsilon_{\alpha\beta}W_{A\alpha L}W_{A\beta L}
  \right]_{\mathcal{F}}\,,
  \end{aligned}
  \tag{27.6.2}\label{eq:27.6.2}
\end{equation}
where \(\mathcal{A}_\lambda\) is a general Lorentz- and gauge-invariant
function; `\(\mathcal{D}\cdots\)' denotes terms involving superderivatives or
spacetime derivatives of the preceding arguments; and \(g_\lambda\) is the
\emph{one-loop} effective gauge coupling, given by the same formula as the
one-loop renormalized gauge coupling constant
\begin{equation}
  g_\lambda^{-2}=\text{constant}-2b\ln\lambda\,,
  \tag{27.6.3}\label{eq:27.6.3}
\end{equation}
where \(b\) is the coefficient of \(g^3\) in the Gell-Mann--Low function
\(\beta(g)\), discussed in Chapter 18. This is for a simple gauge group, with
a single gauge coupling, but the extension to a direct product of simple and
\(U(1)\) gauge groups is trivial. Note in particular that the effective
superpotential is not only finite in the limit \(\lambda\longrightarrow
\infty\), but at least in perturbation theory it contains no terms beyond
those in the original superpotential, and there is no change in the
coefficients of the terms it does contain.

To prove this theorem, we shall interpret this theory as a special case of
one with two additional external gauge-invariant left-chiral superfields
\(X\) and \(Y\), with Lagrangian density
\begin{equation}
  \mathcal{L}^{\#}
  =
  \left[\Phi^\dagger e^{-V}\Phi\right]_D
  +2\operatorname{Re}\left[Yf(\Phi)\right]_{\mathcal{F}}
  +\frac12\operatorname{Re}
  \left[
    X\sum_{A\alpha\beta}
    \epsilon_{\alpha\beta}W_{A\alpha L}W_{A\beta L}
  \right]_{\mathcal{F}}\,.
  \tag{27.6.4}\label{eq:27.6.4}
\end{equation}
This Lagrangian density becomes equal to the original one when the scalar
components \(x\) and \(y\) of \(X\) and \(Y\) are given the values
\(x=1/g^2\) and \(y=1\), and the spinor and auxiliary components of \(X\)
and \(Y\) are set equal to zero. Since supersymmetry and gauge invariance are
assumed to be preserved in the cut-off procedure, the effective Lagrangian
density in the presence of these external superfields must be the sum of the
\(D\)-term of a general superfield and the real part of the
\(\mathcal{F}\)-term of a left-chiral superfield:
\begin{equation}
  \begin{aligned}
  \mathcal{L}^{\#}_\lambda
  ={}&
  \left[
    \mathcal{A}_\lambda
    \bigl(\Phi,\Phi^\dagger,V,X,X^\dagger,Y,Y^\dagger,
    \mathcal{D}\cdots\bigr)
  \right]_D
  \\
  &+2\operatorname{Re}
  \left[
    \mathcal{B}_\lambda(\Phi,W_L,X,Y)
  \right]_{\mathcal{F}}\,,
  \end{aligned}
  \tag{27.6.5}\label{eq:27.6.5}
\end{equation}
with \(\mathcal{A}_\lambda\) and \(\mathcal{B}_\lambda\) both
gauge-invariant functions of the displayed arguments. We do not include any
superderivatives or spacetime derivatives in the \(\mathcal{F}\)-term
because, as in Section 26.3, terms involving derivatives of any of the
left-chiral superfields or their adjoints may be rewritten as contributions
to \([\mathcal{A}_\lambda]_D\). (It is true that Eq.~\eqref{eq:27.3.12}
gives \(W_L\) itself in terms of two \(\mathcal{D}_R\)s acting on a
superfield \(\exp(-2V)\mathcal{D}_L\exp(2V)\), but this superfield is not
gauge-invariant, and we are requiring that \(\mathcal{A}_\lambda\) be
gauge-invariant.)

The dependence of \(\mathcal{B}_\lambda\) on \(X\) and \(Y\) is severely
limited by two additional symmetries of the action obtained from the
Lagrangian density \eqref{eq:27.6.4}. (Both of these symmetries are broken by
non-perturbative effects, which will be considered in Chapter 29.) The first
symmetry is a perturbative \(U(1)\) \(R\)-symmetry, of the sort discussed in
Section 26.3, for which \(\theta_L\) and \(\theta_R\) are given the \(R\)
values \(+1\) and \(-1\), the superfields \(\Phi\), \(V\), and \(X\) are
\(R\)-neutral, and \(Y\) has the \(R\) value \(+2\). (Recall that
\(f_{\mathcal{F}}\) is the coefficient of \(\theta_L^2\) in \(f\), so in
order for \(f_{\mathcal{F}}\) to have the \(R\) value 0, \(f\) must have
the \(R\) value 2.) Because \(W_L\) is given by two \(\mathcal{D}_R\)s and
one \(\mathcal{D}_L\) acting on \(R\)-neutral superfields, it has the \(R\)
value \(+1\). Now, \(R\) invariance requires \(\mathcal{B}_\lambda\), like
the superpotential, to have the \(R\) value \(+2\). It cannot depend on any
superfields with negative \(R\) values, such as adjoints of left-chiral
superfields, because it is \emph{holomorphic}, so the terms in
\(\mathcal{B}_\lambda\) can only be of first order in \(Y\) or of second
order in \(W_L\), with coefficients depending only on the \(R\)-neutral
superfields \(\Phi\) and/or \(X\):
\begin{equation}
  \begin{aligned}
  \mathcal{B}_\lambda(\Phi,W_L,X,Y)
  ={}&
  Yf_\lambda(\Phi,X)
  \\
  &+\sum_{\alpha\beta AB}
  \epsilon_{\alpha\beta}W_{A\alpha L}W_{B\beta L}
  h_{\lambda AB}(\Phi,X)\,.
  \end{aligned}
  \tag{27.6.6}\label{eq:27.6.6}
\end{equation}
(Lorentz invariance requires the spinor indices on the \(W_L\)s to be
contracted with \(\epsilon_{\alpha\beta}\).) The other symmetry is
translation of \(X\) by an imaginary numerical constant,
\(X\longrightarrow X+i\xi\), with \(\xi\) real. This changes the Lagrangian
density \eqref{eq:27.6.4} by an amount proportional to
\(\operatorname{Im}\sum_{A\alpha\beta}W_{A\alpha L}W_{A\beta L}\), which as
we saw in Section 27.3 is a spacetime derivative, and therefore can have no
effect in perturbation theory. This translation symmetry prevents \(X\)
from appearing anywhere in the effective Lagrangian density
\eqref{eq:27.6.5} except where it appears in the original Lagrangian density
\eqref{eq:27.6.4}. We therefore conclude that \(f_\lambda\) is independent
of \(X\), while \(h_{\lambda AB}\) consists of a \(\Phi\)-independent term
proportional to \(X\delta_{AB}\), plus a term that is independent of \(X\).
That is,
\begin{equation}
  \begin{aligned}
  \mathcal{B}_\lambda(\Phi,W_L,X,Y)
  ={}&
  Yf_\lambda(\Phi)
  \\
  &+\sum_{\alpha\beta AB}
  \epsilon_{\alpha\beta}W_{A\alpha L}W_{B\beta L}
  \left[
    c_\lambda\delta_{AB}X+\ell_{\lambda AB}(\Phi)
  \right]\,,
  \end{aligned}
  \tag{27.6.7}\label{eq:27.6.7}
\end{equation}
where \(c_\lambda\) is a real cut-off-dependent constant.

The point of introducing the external auxiliary superfields \(X\) and \(Y\)
is that, by giving them suitable values, we can make use of weak-coupling
approximations to determine the coefficients in Eq.~\eqref{eq:27.6.7}. If we
set the spinor and auxiliary components of \(X\) and \(Y\) equal to zero,
and take their scalar components \(x\) and \(y\) to approach infinity and
zero, respectively, then the gauge coupling constant vanishes as
\(1/\sqrt{x}\), and all Yukawa and scalar couplings derived from the
superpotential vanish as \(y\). In this limit, the only graph that contributes
to the term in \eqref{eq:27.6.7} proportional to \(Y\) has a single vertex
arising from the term \(2\operatorname{Re}[Yf(\Phi)]_{\mathcal{F}}\) in
Eq.~\eqref{eq:27.6.4}, so
\begin{equation}
  f_\lambda(\Phi)=f(\Phi)\,.
  \tag{27.6.8}\label{eq:27.6.8}
\end{equation}
Also, with \(Y=0\) there is a conservation law which requires every term in
\(\mathcal{L}^{\#}_\lambda\) to have equal numbers of \(\Phi\)s and
\(\Phi^\dagger\)s, so since \(\Phi^\dagger\) cannot occur in
\(\ell_{\lambda AB}\), neither can \(\Phi\). Gauge invariance then requires
the constant \(\ell_{\lambda AB}\) to be proportional to \(\delta_{AB}\)
for a simple group:
\begin{equation}
  \ell_{\lambda AB}=\delta_{AB}L_\lambda\,.
  \tag{27.6.9}\label{eq:27.6.9}
\end{equation}
Now, since gauge propagators go as \(1/x\) while pure gauge interactions go
as \(x\) and scalar propagators and interactions are \(x\)-independent, with
\(y=0\) the number of powers of \(x\) in a diagram with \(V_W\) pure gauge
boson vertices, \(I_W\) internal gauge boson lines, and any number of
scalar-gauge boson vertices and scalar propagators is
\begin{equation}
  N_x=V_W-I_W\,.
  \tag{27.6.10}\label{eq:27.6.10}
\end{equation}
The number of loops is given by
\begin{equation}
  L=I_W+I_\Phi-V_W-V_\Phi+1\,,
  \tag{27.6.11}\label{eq:27.6.11}
\end{equation}
where \(I_\Phi\) is the number of internal \(\Phi\) lines and \(V_\Phi\) is
the number of \(\Phi\)-\(V\) interaction vertices. All the \(\Phi\)-\(V\)
vertices have two \(\Phi\) lines attached, so with no external \(\Phi\)
lines \(I_\Phi\) and \(V_\Phi\) are equal, and therefore cancel in
Eq.~\eqref{eq:27.6.11}, so that Eq.~\eqref{eq:27.6.10} may be written
\begin{equation}
  N_x=1-L\,.
  \tag{27.6.12}\label{eq:27.6.12}
\end{equation}
Thus the coefficient \(c_\lambda\) of \(X\) in Eq.~\eqref{eq:27.6.7} is
correctly given by the tree approximation and is therefore what it was in
the original Lagrangian, simply \(c_\lambda=1\), while the coefficient
\(L_\lambda\) of the \(X\)-independent term is given by one-loop diagrams
only. Putting this all together, we have
\begin{equation}
  \begin{aligned}
  \mathcal{L}^{\#}_\lambda
  ={}&
  \left[
    \mathcal{A}_\lambda
    \bigl(\Phi,\Phi^\dagger,V,X,X^\dagger,Y,Y^\dagger,
    \mathcal{D}\cdots\bigr)
  \right]_D
  +2\operatorname{Re}\left[Yf(\Phi)\right]_{\mathcal{F}}
  \\
  &+\frac12\operatorname{Re}
  \left[
    (X+L_\lambda)
    \sum_{A\alpha\beta}
    \epsilon_{\alpha\beta}W_{A\alpha L}W_{A\beta L}
  \right]_{\mathcal{F}}\,,
  \end{aligned}
  \tag{27.6.13}\label{eq:27.6.13}
\end{equation}
where \(L_\lambda\) is the one-loop contribution. Setting \(Y=1\) and
\(X=1/g^2\) then gives Eq.~\eqref{eq:27.6.2}, with
\(g_\lambda^{-2}=g^{-2}+L_\lambda\). As shown in Section 18.3, the leading
order contribution to \(\lambda\,dg_\lambda/d\lambda\) is the same function
of \(g_\lambda\) whatever renormalization scheme is used to define this
coupling, so to one-loop order we must have
\begin{equation}
  \lambda\,\frac{dg_\lambda}{d\lambda}=bg_\lambda^3\,,
  \tag{27.6.14}\label{eq:27.6.14}
\end{equation}
where \(b\) is the same coefficient of \(g^3\) as in the renormalization
group equation of Gell-Mann and Low. The solution is
Eq.~\eqref{eq:27.6.3}, completing the proof.

In theories with a \(U(1)\) gauge superfield \(V_1\), the Lagrangian may
contain a Fayet--Iliopoulos term \eqref{eq:27.2.7}:
\begin{equation}
  \mathcal{L}_{\mathrm{FI}}=\xi[V_1]_D\,.
  \tag{27.6.15}\label{eq:27.6.15}
\end{equation}
It is easy to see that the coefficient \(\xi\) of such a term is not
renormalized~\hyperref[ch27-ref-7]{[7]}. If the corresponding coefficient
\(\xi_\lambda\) in the Wilsonian Lagrangian density did depend on the gauge
couplings or the couplings in the superpotential, then when we replace the
original Lagrangian \eqref{eq:27.6.1} with a Lagrangian
\eqref{eq:27.6.4} involving the external superfields \(X\) and \(Y\), this
term in the Wilsonian Lagrangian would be required by supersymmetry to take
the form
\begin{equation}
  \mathcal{L}^{\#}_{\mathrm{FI}\,\lambda}
  =
  \left[
    \xi_\lambda(X,Y,X^*,Y^*)V_1
  \right]_D\,,
  \tag{27.6.16}\label{eq:27.6.16}
\end{equation}
with \(\xi_\lambda\) a function with a non-trivial dependence on \(X\)
and/or \(Y\) and/or their adjoints. But such a term would not be
gauge-invariant, because, according to Eq.~\eqref{eq:27.2.18}, a gauge
transformation shifts \(V_1\) by a chiral superfield
\(i(\Omega-\Omega^*)/2\), and although the \(D\)-term of a chiral superfield
vanishes, the product of \(i(\Omega-\Omega^*)/2\) and \(\xi_\lambda\) is not
chiral for general gauge transformations if \(\xi_\lambda\) has any
dependence on other superfields. There actually are diagrams that make
contributions to \(\xi_\lambda\) that are independent

\IfFileExists{chapters/chapter27/figures/fig27-01.tex}{%
  % Source: physical PDF p. 176 (printed p. 153). Coverage: complete Figure
% 27.1 and caption. Uncertainty: none; line styles, loop orientation, and
% arrow directions reconstructed from the full-resolution rendered scan.
\begin{figure}[H]
  \centering
  \begin{tikzpicture}[
    x=1cm,
    y=1cm,
    line cap=round,
    line join=round,
    scalar/.style={
      draw,
      line width=0.8pt,
      dash pattern=on 4.2pt off 4.2pt
    },
    flow/.style={
      draw,
      line width=0.8pt,
      -{Stealth[length=2.2mm,width=1.7mm]}
    }
  ]
    % A single incoming complex-scalar line enters a clockwise oriented
    % complex-scalar loop. All lines are dashed in the source.
    \draw[scalar] (-2.75,0) -- (0.45,0);
    \draw[flow] (-1.95,0) -- (-1.42,0);

    \draw[scalar] (1.40,0) circle[radius=0.95];
    \draw[flow] (1.18,0.95) -- (1.67,0.95);
    \draw[flow] (1.63,-0.95) -- (1.14,-0.95);
  \end{tikzpicture}
  \renewcommand{\thefigure}{27.1}
  \caption{A one-loop diagram that could be quadratically divergent in
  theories with supersymmetry broken by trilinear couplings among scalar
  fields and their adjoints. The lines all represent complex scalar fields.}
  \label{fig:27.1}
\end{figure}
%
}{%
  % Source: physical PDF p. 176 (printed p. 153). Coverage: complete Figure
% 27.1 and caption. Uncertainty: none; line styles, loop orientation, and
% arrow directions reconstructed from the full-resolution rendered scan.
\begin{figure}[H]
  \centering
  \begin{tikzpicture}[
    x=1cm,
    y=1cm,
    line cap=round,
    line join=round,
    scalar/.style={
      draw,
      line width=0.8pt,
      dash pattern=on 4.2pt off 4.2pt
    },
    flow/.style={
      draw,
      line width=0.8pt,
      -{Stealth[length=2.2mm,width=1.7mm]}
    }
  ]
    % A single incoming complex-scalar line enters a clockwise oriented
    % complex-scalar loop. All lines are dashed in the source.
    \draw[scalar] (-2.75,0) -- (0.45,0);
    \draw[flow] (-1.95,0) -- (-1.42,0);

    \draw[scalar] (1.40,0) circle[radius=0.95];
    \draw[flow] (1.18,0.95) -- (1.67,0.95);
    \draw[flow] (1.63,-0.95) -- (1.14,-0.95);
  \end{tikzpicture}
  \renewcommand{\thefigure}{27.1}
  \caption{A one-loop diagram that could be quadratically divergent in
  theories with supersymmetry broken by trilinear couplings among scalar
  fields and their adjoints. The lines all represent complex scalar fields.}
  \label{fig:27.1}
\end{figure}
%
}

of all coupling constants. For the Lagrangian \eqref{eq:27.6.1} there are no
factors of the gauge
coupling \(g\) at vertices at which the gauge superfield interacts with
chiral matter, but instead a factor \(g^{-2}\) for each gauge propagator, so
a graph with no internal gauge lines and no self-couplings of the chiral
superfield will have no dependence on coupling constants. The only such
graphs that contribute to \(\xi_\lambda\) are those in which a single
external gauge line is attached to a chiral loop. (See Figure~\ref{fig:27.1}.)
The contribution of all such graphs is proportional to the sum of the gauge
couplings of all chiral superfields---that is, to the trace of the \(U(1)\)
generator. But as discussed in Section 22.4, this trace must vanish (if the
\(U(1)\) symmetry is unbroken) in order to avoid gravitational anomalies
that violate the conservation of the \(U(1)\) current.

The most important application of these theorems is a corollary, which tells
us that if there is no Fayet--Iliopoulos term and if the superpotential
\(f(\Phi)\) allows solutions of the equations
\(\partial f(\phi)/\partial\phi_n=0\), then supersymmetry is not broken in
any finite order of perturbation theory.

To test this we must examine Lorentz-invariant field configurations, in which
the \(\Phi_n\) have only constant scalar components \(\phi_n\) and constant
auxiliary components \(\mathcal{F}_n\), while (in Wess--Zumino gauge) the
coefficients \(V_A\) of the gauge generators \(t_A\) in the matrix gauge
superfield \(V\) have only auxiliary components \(D_A\). Supersymmetry is
unbroken if there are values of \(\phi_n\) for which
\(\mathcal{L}_\lambda\) has no terms of first order in \(\mathcal{F}_n\) or
\(D_A\), in which case there is sure to be an equilibrium solution with
\(\mathcal{F}_n=D_A=0\). (In Section 29.2 we will see that this is the
sufficient as well as the necessary condition for supersymmetry to be
unbroken.) In the absence of Fayet--Iliopoulos terms, this will be the case
if for all \(A\)
\begin{equation}
  \sum_{nm}
  \frac{\partial K_\lambda(\phi,\phi^*)}{\partial\phi_n^*}
  (t_A)_{mn}\phi_m^*
  =0
  \tag{27.6.17}\label{eq:27.6.17}
\end{equation}
and for all \(n\)
\begin{equation}
  \frac{\partial f(\phi)}{\partial\phi_n}=0\,,
  \tag{27.6.18}\label{eq:27.6.18}
\end{equation}
where the effective K\"ahler potential \(K_\lambda(\phi,\phi^*)\) is
\begin{equation}
  K_\lambda(\phi,\phi^*)
  =
  \mathcal{A}_\lambda(\phi,\phi^*,0,0\cdots)\,,
  \tag{27.6.19}\label{eq:27.6.19}
\end{equation}
with \(\mathcal{A}_\lambda(\phi,\phi^*,0,0\cdots)\) obtained from
\(\mathcal{A}_\lambda\) by setting the gauge superfield and all
superderivatives equal to zero. (With superderivatives required to vanish by
Lorentz invariance, the only dependence of \(\mathcal{A}_\lambda\) on \(V\)
is a factor \(\exp(-V)\) following every factor \(\Phi^\dagger\).)

We now use a trick that we have already employed in Section 27.4. If there is
any solution \(\phi^{(0)}\) of Eq.~\eqref{eq:27.6.18}, then the gauge
symmetry tells us that there is a continuum of such solutions, with
\(\phi_n\) replaced with
\begin{equation}
  \phi_n(z)
  =
  \left[
    \exp\left(i\sum_A t_Az_A\right)
  \right]_{nm}
  \phi_m^{(0)}\,,
  \tag{27.6.20}\label{eq:27.6.20}
\end{equation}
where (since \(f\) depends only on \(\phi\), not \(\phi^*\)) the \(z_A\)
are an arbitrary set of \emph{complex} parameters. If
\(K_\lambda(\phi,\phi^*)\) has a stationary point anywhere on the surface
\(\phi=\phi(z)\), then at that point
\begin{equation}
  \begin{aligned}
  0={}&
  \sum_{nmA}
  \frac{\partial K_\lambda(\phi,\phi^*)}{\partial\phi_n}
  (t_A)_{nm}\phi_m\,\delta z_A
  \\
  &-
  \sum_{nmA}
  \frac{\partial K_\lambda(\phi,\phi^*)}{\partial\phi_n^*}
  (t_A)_{mn}\phi_m^*\,\delta z_A^*\,.
  \end{aligned}
  \tag{27.6.21}\label{eq:27.6.21}
\end{equation}
Since this must be satisfied for all infinitesimal complex \(\delta z_A\),
the coefficients of both \(\delta z_A\) and \(\delta z_A^*\) must both
vanish, and therefore Eq.~\eqref{eq:27.6.17} as well as
Eq.~\eqref{eq:27.6.18} is satisfied at this point. Thus the existence of a
stationary point of \(K_\lambda(\phi,\phi^*)\) on the surface
\(\phi=\phi(z)\) would imply that supersymmetry is unbroken to all orders of
perturbation theory.

The zeroth-order K\"ahler potential \((\phi^\dagger\phi)\) is bounded below
and goes to infinity as \(\phi\longrightarrow\infty\), so it certainly has a
minimum on the surface \(\phi=\phi(z)\), where of course it is stationary.
If there were no flat directions in which \(K_\lambda\) is constant at this
minimum then any sufficiently small perturbation to the K\"ahler potential
might shift the minimum, but would not destroy it. At the minimum of the
K\"ahler potential on the surface \(\phi=\phi(z)\) there are flat
directions: ordinary global gauge transformations
\(\delta\phi=i\sum_A\delta z_A t_A\phi\) with \(z_A\) real. But these are
also flat directions for the perturbation
\(K_\lambda(\phi,\phi^*)-(\phi^\dagger\phi)\), so there is still a local
minimum of \(K_\lambda\) on the surface \(\phi=\phi(z)\) for any
perturbation in at least a finite range, and thus to all orders in whatever
couplings appear in \(K_\lambda(\phi,\phi^*)\). As we have seen, this is a
set of scalar field values at which \(\mathcal{F}_n=0\) and \(D_A=0\) for
all \(n\) and \(A\), which means that supersymmetry is unbroken.

\begin{center}
\(\ast\qquad\ast\qquad\ast\)
\end{center}

These results may be extended to non-renormalizable
theories~\hyperref[ch27-ref-3]{[3]}. In such theories the first term
\([\Phi^\dagger e^{-V}\Phi]_D\) in Eq.~\eqref{eq:27.6.1} is replaced with
the \(D\)-term of an arbitrary real gauge-invariant scalar function of
\(\Phi^\dagger\), \(\Phi\), \(V\), and their superderivatives and spacetime
derivatives, while the second and third terms in Eq.~\eqref{eq:27.6.1} are
replaced with the \(\mathcal{F}\)-term of an arbitrary globally
gauge-invariant scalar function \(f(\Phi,W)\) of \(\Phi_n\) and \(W_\alpha\).
It has been shown that, to all orders of perturbation theory, the function
\(f_\lambda(\Phi,W)\) appearing in the \(\mathcal{F}\)-term of the
Wilsonian Lagrangian is the same as \(f(\Phi,W)\), except for the one-loop
renormalization of the term quadratic in \(W\).
